\chapter*{Pendahuluan} \label{intro:chapter}
%mbxSTARTIGNORE
\addfakecontentsline{Pendahuluan}
\markboth{PENDAHULUAN}{PENDAHULUAN}
%mbxENDIGNORE

%%%%%%%%%%%%%%%%%%%%%%%%%%%%%%%%%%%%%%%%%%%%%%%%%%%%%%%%%%%%%%%%%%%%%%%%%%%%%%
\section{Catatan tentang catatan ini}
\label{notes:section}

%mbxINTROSUBSECTION

\sectionnotes{Bagian untuk pengajar.}

Buku ini berawal dari catatan kuliah saya untuk Math 286
di \href{https://www.math.uiuc.edu/}{University of Illinois di
Urbana-Champaign} (UIUC)
pada Musim Gugur 2008 dan
Musim Semi 2009.
Ini adalah mata kuliah pertama tentang persamaan diferensial untuk para insinyur.
Dengan menggunakan buku ini, saya juga mengajar Math 285 di UIUC\@,
Math 20D di
\href{https://www.math.ucsd.edu/}{University of California, San Diego} (UCSD),
serta Math 2233 dan 4233 di
\href{https://math.okstate.edu/}{Oklahoma State University} (OSU).
Biasanya mata kuliah-mata kuliah ini diajarkan dengan menggunakan
Edwards dan Penney, \emph{Differential
Equations and Boundary Value Problems: Computing and Modeling}~\cite{EP}, atau
Boyce dan DiPrima,
\emph{Elementary
Differential Equations and Boundary Value Problems}~\cite{BD},
dan buku ini dimaksudkan sebagai pengganti yang kurang lebih dapat langsung digunakan.
Buku-buku lain yang saya gunakan sebagai sumber informasi dan inspirasi
adalah karya klasik (dan murah) E.L.\ Ince,
\emph{Ordinary Differential Equations}~\cite{I},
Stanley Farlow, \emph{Differential Equations and Their
Applications}~\cite{F}, yang kini tersedia dari Dover,
Berg dan McGregor,
\emph{Elementary Partial Differential Equations}~\cite{BM},
serta buku gratis karya William Trench,
\emph{Elementary
Differential Equations with Boundary Value Problems}~\cite{T}.
Lihat bab \hyperref[furtherreading:chapter]{Bacaan Lanjutan} di bagian akhir buku.

\subsection{Organisasi}

Pengorganisasian buku ini, sampai tingkat tertentu,
mengharuskan bab-babnya dibahas secara berurutan.
Bab-bab selanjutnya dapat dilewati.
Ketergantungan materi yang dibahas kira-kira adalah sebagai berikut:

% If changing make sure to also update figures/chapterdiagram.pdf_t
% That's a hopefully short term hack before I figure out how to do it
% better so that it also gets links and such
%mbxSTARTIGNORE
\begin{equation*}
\begin{tikzcd}[cramped, row sep=small]
& {\text{\hyperref[intro:chapter]{Pendahuluan}}} \arrow[d] \\
{\text{\Appendixref{linalg:appendix}}} \arrow[dd, dotted]
& {\text{\Chapterref{fo:chapter}}} \arrow[d] \\
& {\text{\Chapterref{ho:chapter}}} \arrow[dddr] \arrow[dd] \arrow[dl] \arrow[dr] \\
{\text{\Chapterref{sys:chapter}}} \arrow[dr, dotted] \arrow[d] & &
  {\text{\Chapterref{ps:chapter}}} \\
{\text{\Chapterref{nlin:chapter}}} & {\text{\Chapterref{FS:chapter}}} \arrow[d]
\arrow[dr,dotted] \\
& {\text{\Chapterref{SL:chapter}}}
& {\text{\Chapterref{LT:chapter}}}
\end{tikzcd}
\end{equation*}
%mbxENDIGNORE
%mbxlatex \begin{center}
%mbxlatex \inputpdft{chapterdiagram}{%
%mbxlatex Diagram dengan 'Pendahuluan' di bagian atas, dengan sebuah panah
%mbxlatex menuju 'Bab 1', lalu sebuah panah menuju 'Bab 2', kemudian empat
%mbxlatex panah menuju 'Bab 3', 'Bab 4', 'Bab 6', dan
%mbxlatex 'Bab 7'. 'Bab 3' memiliki sebuah panah menuju 'Bab 8'
%mbxlatex dan sebuah panah putus-putus menuju 'Bab 4'. 'Bab 4' memiliki
%mbxlatex sebuah panah menuju 'Bab 5' dan sebuah panah putus-putus menuju 'Bab 6'.
%mbxlatex Ada juga 'Lampiran A' dengan sebuah panah putus-putus menuju 'Bab 3'.}
%mbxlatex \end{center}

Ada beberapa rujukan dalam bab \ref{FS:chapter} dan \ref{SL:chapter}
ke \chapterref{sys:chapter} (sebagian aljabar linear), tetapi rujukan tersebut
tidak esensial dan dapat dibaca sekilas,
sehingga \chapterref{sys:chapter}
dapat dengan aman dilewati, sementara
bab \ref{FS:chapter} dan \ref{SL:chapter} tetap dibahas.
\Chapterref{LT:chapter} tidak bergantung pada
\chapterref{FS:chapter}, kecuali bahwa
bagian PDE pada \ref{laplacepde:section} membuat
beberapa rujukan ke
\chapterref{FS:chapter},
meskipun secara teori bagian tersebut dapat dibahas
secara terpisah.
\appendixref{linalg:appendix} yang lebih mendalam tentang aljabar linear
dapat menggantikan ulasan singkat pada \sectionref{sec:matrix}
untuk mata kuliah yang menggabungkan aljabar linear dan ODE.

\subsection{Mata kuliah yang umum}

Buku ini dapat digunakan untuk menyelenggarakan beberapa mata kuliah yang berbeda.
Ada dua mata kuliah satu semester yang semula diajarkan di UIUC\@.
Keduanya mencakup ODE dan juga beberapa PDE.
Mata kuliah pertama berlangsung selama 4 jam per minggu (Math 286):

\medskip

\noindent
\hyperref[intro:chapter]{Pendahuluan} (\ref{introde:section}),
\chapterref{fo:chapter} (\ref{integralsols:section}--\ref{numer:section}),
\chapterref{ho:chapter},
\chapterref{sys:chapter},
\chapterref{FS:chapter} (\ref{bvp:section}--\ref{dirich:section}),
\chapterref{SL:chapter} (atau
\ref{LT:chapter} atau \ref{ps:chapter} atau \ref{nlin:chapter}).

\medskip

Mata kuliah kedua di UIUC berlangsung 3 jam per minggu (Math 285):

\medskip

\noindent
\hyperref[intro:chapter]{Pendahuluan} (\ref{introde:section}),
\chapterref{fo:chapter} (\ref{integralsols:section}--\ref{numer:section}),
\chapterref{ho:chapter},
\chapterref{FS:chapter} (\ref{bvp:section}--\ref{dirich:section}),
dan mungkin \chapterref{SL:chapter},
\ref{LT:chapter}, atau \ref{ps:chapter}.

\medskip

Mata kuliah satu semester dengan 3 jam per minggu yang tidak membahas sistem maupun PDE
akan mencakup, selain pendahuluan,
%\sectionref{introde:section},
\chapterref{fo:chapter},
\chapterref{ho:chapter},
\chapterref{LT:chapter}, dan \chapterref{ps:chapter},
(dengan beberapa bagian dilewati seperti di atas).
Sebaliknya, mata kuliah umum yang membahas sistem mungkin perlu melewati transformasi Laplace dan deret pangkat,
lalu membahas
%\sectionref{introde:section},
\chapterref{fo:chapter},
\chapterref{ho:chapter},
\chapterref{sys:chapter}, dan \chapterref{nlin:chapter}.

\medskip

Jika beberapa bagian perlu dilewati pada awalnya, inti yang baik dari
bagian-bagian tentang ODE tunggal adalah:
\ref{introde:section},
\ref{integralsols:section}--\ref{intfactor:section},
\ref{auteq:section},
\ref{solinear:section},
\ref{sec:ccsol},
\ref{sec:mv}--\ref{forcedo:section}.

\medskip

Buku lengkap dapat dibahas dengan laju yang cukup cepat dalam kira-kira 76 pertemuan
(tanpa \appendixref{linalg:appendix})
atau 86 pertemuan (dengan \appendixref{linalg:appendix} menggantikan
\sectionref{sec:matrix}).
Ini belum memperhitungkan ujian, peninjauan ulang,
atau waktu yang dihabiskan di laboratorium komputer. % (if using IODE for example).
Dengan materi ini, mata kuliah dua kuartal atau dua semester dapat diselenggarakan dengan mudah.
Sebagai contoh (dengan kemungkinan melewati beberapa bagian secara strategis):

\medskip

\noindent
Semester 1:
\hyperref[intro:chapter]{Pendahuluan},
\chapterref{fo:chapter},
\chapterref{ho:chapter},
\chapterref{LT:chapter},
\chapterref{ps:chapter}.
\\
Semester 2:
\Chapterref{sys:chapter},
\chapterref{nlin:chapter},
\chapterref{FS:chapter},
\chapterref{SL:chapter}.

\medskip

Mata kuliah gabungan tentang ODE dan aljabar linear dapat disusun sebagai berikut:

\medskip

\noindent
\hyperref[intro:chapter]{Pendahuluan},
\chapterref{fo:chapter} (\ref{integralsols:section}--\ref{numer:section}),
\chapterref{ho:chapter},
\appendixref{linalg:appendix},
\chapterref{sys:chapter} (tanpa \sectionref{sec:matrix}), mungkin
\chapterref{nlin:chapter}.

\medskip

Bab tentang
transformasi Laplace (\chapterref{LT:chapter}),
bab tentang Sturm--Liouville (\chapterref{SL:chapter}),
bab tentang deret pangkat (\chapterref{ps:chapter}),
dan bab tentang sistem nonlinear (\chapterref{nlin:chapter}),
kurang lebih dapat saling dipertukarkan dan dapat diperlakukan sebagai \myquote{topik}.
Jika \chapterref{nlin:chapter} dibahas, sebaiknya bab itu ditempatkan tepat
setelah \chapterref{sys:chapter},
dan \chapterref{SL:chapter} paling baik dibahas tepat setelah
\chapterref{FS:chapter}.
Jika waktu terbatas, dua bagian pertama dari
\chapterref{ps:chapter} merupakan unit mandiri yang cukup baik.

%\medskip
\subsection{Sumber daya komputer}

Situs web buku ini,
\url{https://www.jirka.org/diffyqs/}
atau
\url{https://jirilebl.github.io/diffyqs/},
memuat sumber daya berikut:
\begin{enumerate}
\item Demo SAGE interaktif.
\item Pekerjaan rumah WeBWorK daring
(dengan menggunakan instalasi WeBWorK sendiri atau Edfinity)
untuk sebagian besar bagian, yang disesuaikan untuk buku ini.
\item PDF gambar-gambar yang digunakan dalam buku ini.
\item Video YouTube dan slide yang sesuai untuk beberapa bagian.
\end{enumerate}

Saya menggunakan IODE\index{perangkat lunak IODE}
(\url{https://publish.illinois.edu/iode-diffeq/})
dalam mata kuliah UIUC.
IODE adalah paket perangkat lunak gratis yang
bekerja dengan Matlab (perangkat lunak proprietari) atau Octave (perangkat lunak gratis).
%Unfortunately IODE is not kept up to date at this point, and may have
%trouble running on newer versions of Matlab.
Grafik dalam buku ini dibuat dengan
perangkat lunak Genius\index{perangkat lunak Genius}
(lihat \url{https://www.jirka.org/genius.html}).
%  I use Genius
%in class to show these (and other) graphs.

Kode sumber \LaTeX\ buku ini juga tersedia
untuk kemungkinan modifikasi dan penyesuaian
di GitHub (\url{https://github.com/jirilebl/diffyqs}).

%\medskip

%\textbf{Ucapan terima kasih:}

\subsection{Ucapan terima kasih}

Pertama-tama, saya ingin menyampaikan terima kasih kepada Rick Laugesen. Saya menggunakan catatan kuliahnya yang ditulis tangan
ketika pertama kali mengajar
Math 286. Pengorganisasian buku ini sampai Bab 5,
serta pilihan
materi yang dibahas, sangat dipengaruhi oleh catatannya. Banyak
contoh dan perhitungan diambil dari catatannya. Saya juga sangat berutang budi kepada Rick atas semua nasihat yang diberikannya, bukan hanya tentang pengajaran
Math 286.
Atas bantuan menemukan kesalahan dan saran-saran lainnya,
saya juga ingin menyampaikan terima kasih (tanpa urutan tertentu):
John P.\ D'Angelo,
Sean Raleigh, Jessica Robinson, Michael Angelini, Leonardo Gomes, Jeff
Winegar, Ian Simon, Thomas Wicklund, Eliot Brenner, Sean Robinson,
Jannett Susberry, Dana Al-Quadi, Cesar Alvarez, Cem Bagdatlioglu,
Nathan Wong, Alison Shive, Shawn White, Wing Yip Ho, Joanne Shin,
Gladys Cruz, Jonathan Gomez, Janelle Louie, Navid Froutan,
Grace Victorine, Paul Pearson, Jared Teague, Ziad Adwan,
Martin Weilandt, S\"{o}nmez \c{S}ahuto\u{g}lu,
Pete Peterson, Thomas Gresham, Prentiss Hyde, Jai Welch,
Simon Tse, Andrew Browning, James Choi, Dusty Grundmeier,
John Marriott,
Jim Kruidenier,
Barry Conrad,
Wesley Snider,
Colton Koop,
Sarah Morse,
Erik Boczko,
Asif Shakeel,
Chris Peterson,
Nicholas Hu,
Paul Seeburger,
Jonathan McCormick,
David Leep,
William Meisel,
Shishir Agrawal,
Tom Wan,
Andres Valloud,
Martin Irungu,
Justin Corvino,
Tai-Peng Tsai,
Santiago Mendoza Reyes,
Glen Pugh,
Michael Tran,
Heber Farnsworth,
Tam\'as Zsoldos,
Mark Mills,
George Ballinger,
dan mungkin orang-orang lain yang saya lupakan.
Terakhir, saya ingin menyampaikan terima kasih atas hibah NSF DMS-0900885 dan DMS-1362337.


%%%%%%%%%%%%%%%%%%%%%%%%%%%%%%%%%%%%%%%%%%%%%%%%%%%%%%%%%%%%%%%%%%%%%%%%%%%%%%
\sectionnewpage
\section{Pengantar persamaan diferensial}
\label{introde:section}

%mbxINTROSUBSECTION

\sectionnotes{lebih dari 1 kuliah\EPref{, \S1.1 dalam \cite{EP}}\BDref{, bab 1 dalam \cite{BD}}}

\subsection{Persamaan diferensial}

Hukum-hukum fisika umumnya dituliskan sebagai persamaan diferensial.
Oleh karena itu, semua bidang sains dan teknik menggunakan persamaan diferensial sampai tingkat tertentu.
Persamaan diferensial penting untuk memahami hampir segala sesuatu yang akan
Anda pelajari dalam kelas sains dan teknik.
Anda dapat menganggap matematika sebagai bahasa sains, dan
persamaan diferensial adalah salah satu bagian terpenting dari bahasa ini
sejauh menyangkut sains dan teknik. Sebagai analogi,
misalkan mulai sekarang semua kelas Anda diajarkan dalam bahasa Swahili.
Penting untuk mempelajari bahasa Swahili terlebih dahulu, atau Anda akan sangat
kesulitan mendapatkan nilai baik di kelas.

Sebuah \emph{\myindex{persamaan diferensial}} hanyalah persamaan yang melibatkan
bukan hanya variabel tak diketahui, tetapi juga turunannya.
Dalam kalkulus, Anda mungkin telah melihat dan bahkan menyelesaikan banyak
persamaan diferensial tanpa menyadarinya.
Berikut contoh yang mungkin belum pernah Anda lihat:
\begin{equation} \label{eq1}
\frac{dx}{dt} + x = 2 \cos t .
\end{equation}
Besaran tak diketahui, $x$, adalah \emph{\myindex{variabel dependen}} dan $t$ adalah
\emph{\myindex{variabel independen}}. Artinya, $x$ adalah fungsi dari $t$.
Persamaan \eqref{eq1}
adalah contoh \emph{\myindex{persamaan diferensial orde pertama}}, karena
hanya melibatkan turunan pertama variabel dependen.
Persamaan ini muncul dari hukum pendinginan Newton, ketika suhu lingkungan
berosilasi terhadap waktu.
